\documentclass[10pt,a4paper]{amsart}
\usepackage{amsmath,amssymb,amsthm,graphicx}
% Retained minimal preamble for the Stacks Project chapter flat.tex (More on Flatness).
% Verbatim excerpt of the chapter preamble: Xy-pic setup (native preamble lines 14-18) and the native
% theorem declarations with the mathematical macros (native preamble lines 154-217). Copied mechanically;
% nothing added, removed or redefined. The wrapper class and the amsmath/amssymb/amsthm/graphicx base
% packages are supplied by the assembler, so the class/externaldocument part of the native preamble is omitted.
% For commutative diagrams we use Xy-pic
\usepackage[all]{xy}

% We use 2cell for 2-commutative diagrams.
\xyoption{2cell}
% Theorem environments.
%
\theoremstyle{plain}
\newtheorem{theorem}[subsection]{Theorem}
\newtheorem{proposition}[subsection]{Proposition}
\newtheorem{lemma}[subsection]{Lemma}

\theoremstyle{definition}
\newtheorem{definition}[subsection]{Definition}
\newtheorem{example}[subsection]{Example}
\newtheorem{exercise}[subsection]{Exercise}
\newtheorem{situation}[subsection]{Situation}

\theoremstyle{remark}
\newtheorem{remark}[subsection]{Remark}
\newtheorem{remarks}[subsection]{Remarks}

\numberwithin{equation}{subsection}

% Macros
%
\def\lim{\mathop{\mathrm{lim}}\nolimits}
\def\colim{\mathop{\mathrm{colim}}\nolimits}
\def\Spec{\mathop{\mathrm{Spec}}}
\def\Hom{\mathop{\mathrm{Hom}}\nolimits}
\def\Ext{\mathop{\mathrm{Ext}}\nolimits}
\def\SheafHom{\mathop{\mathcal{H}\!\mathit{om}}\nolimits}
\def\SheafExt{\mathop{\mathcal{E}\!\mathit{xt}}\nolimits}
\def\Sch{\mathit{Sch}}
\def\Mor{\mathop{\mathrm{Mor}}\nolimits}
\def\Ob{\mathop{\mathrm{Ob}}\nolimits}
\def\Sh{\mathop{\mathit{Sh}}\nolimits}
\def\NL{\mathop{N\!L}\nolimits}
\def\CH{\mathop{\mathrm{CH}}\nolimits}
\def\proetale{{pro\text{-}\acute{e}tale}}
\def\etale{{\acute{e}tale}}
\def\QCoh{\mathit{QCoh}}
\def\Ker{\mathop{\mathrm{Ker}}}
\def\Im{\mathop{\mathrm{Im}}}
\def\Coker{\mathop{\mathrm{Coker}}}
\def\Coim{\mathop{\mathrm{Coim}}}

% Boxtimes
%
\DeclareMathSymbol{\boxtimes}{\mathbin}{AMSa}{"02}

%
% Macros for moduli stacks/spaces
%
\def\QCohstack{\mathcal{QC}\!\mathit{oh}}
\def\Cohstack{\mathcal{C}\!\mathit{oh}}
\def\Spacesstack{\mathcal{S}\!\mathit{paces}}
\def\Quotfunctor{\mathrm{Quot}}
\def\Hilbfunctor{\mathrm{Hilb}}
\def\Curvesstack{\mathcal{C}\!\mathit{urves}}
\def\Polarizedstack{\mathcal{P}\!\mathit{olarized}}
\def\Complexesstack{\mathcal{C}\!\mathit{omplexes}}
% \Pic is the operator that assigns to X its picard group, usage \Pic(X)
% \Picardstack_{X/B} denotes the Picard stack of X over B
% \Picardfunctor_{X/B} denotes the Picard functor of X over B
\def\Pic{\mathop{\mathrm{Pic}}\nolimits}
\def\Picardstack{\mathcal{P}\!\mathit{ic}}
\def\Picardfunctor{\mathrm{Pic}}
\def\Deformationcategory{\mathcal{D}\!\mathit{ef}}

% --- appended 2026-10-12 (Chief): the project own definitions of the reference/slogan/history
% display environments, copied mechanically from the shared Stacks preamble lines 8-12
% (\usepackage{verbatim} plus the three \newenvironment lines, which the chapter preamble keeps
% verbatim from the project preamble). Nothing added, removed or redefined: these environments
% carry no mathematical content and the project itself typesets them as comments. Without these
% three lines a counted statement whose own text uses them fails with Environment slogan undefined.
\usepackage{verbatim}
\newenvironment{reference}{\comment}{\endcomment}
\newenvironment{slogan}{\comment}{\endcomment}
\newenvironment{history}{\comment}{\endcomment}
\usepackage{cleveref}
\crefname{originaldefinition}{original definition}{original definitions}
\Crefname{originaldefinition}{Original definition}{Original definitions}
\makeatletter
\newlabel{topologies-definition-big-small-ph}{{0DBL}{source}{}{sourcepointer-1}{}}
\newlabel{topologies-definition-big-small-ph@cref}{{[originaldefinition][0][]0DBL}{[1][0][]source}{}{}{}}
\makeatother
\crefname{originalequation}{original equation}{original equations}
\Crefname{originalequation}{Original equation}{Original equations}
\makeatletter
\newlabel{equation-blow-up-square}{{0EV5}{source}{}{sourcepointer-2}{}}
\newlabel{equation-blow-up-square@cref}{{[originalequation][0][]0EV5}{[1][0][]source}{}{}{}}
\makeatother
\crefname{originallemma}{original lemma}{original lemmas}
\Crefname{originallemma}{Original lemma}{Original lemmas}
\makeatletter
\newlabel{topologies-lemma-surjective-proper-ph}{{0DES}{source}{}{sourcepointer-3}{}}
\newlabel{topologies-lemma-surjective-proper-ph@cref}{{[originallemma][0][]0DES}{[1][0][]source}{}{}{}}
\makeatother
\title{Every ph sheaf turns a blowup square into a cartesian diagram of sets}
\author{The Stacks Project Authors}
\date{stacks-project:flat.tex:local-snapshot-20261011; GNU FDL 1.2 or later}
\begin{document}
\maketitle
\noindent\textit{Editorial note.} All delivered mathematical text is the author's own, copied from the pinned Stacks Project chapter file flat.tex; no statement or proof step was rewritten, reordered or removed.\par\medskip
\noindent\textit{Editorial note.} The preamble is the chapter's own retained minimal preamble (sources/stacks-flat/retained-preamble.tex), a verbatim excerpt of the native Xy-pic setup and of the native theorem declarations and mathematical macros, followed by the project own definitions of the reference, slogan and history display environments (verbatim package plus three newenvironment lines) that the retained file carries at its end and that the project itself typesets as comments. No macro was invented or redefined, and the wrapper class and the amsmath/amssymb/amsthm/graphicx base packages are supplied by the assembler.\par\medskip
\noindent\textit{Editorial note.} The counted statement and proof contain no cross reference, no citation and no included file, so the delivered body is self contained.\par\medskip
\noindent\textit{Editorial note.} Canonical location of the counted proposition: Stacks tag 0EW1 (sources/stacks-flat/pinned-original-source.tex line 10878).\par\medskip
\noindent\textit{Allowed standard prerequisites (editorial).}\begin{itemize}
\item Surjective proper morphisms are ph coverings and the big ph site. Topologies, Lemma labelled topologies-lemma-surjective-proper-ph, cited at proof lines 10892-10894, with Topologies, Definition labelled topologies-definition-big-small-ph invoked in the statement. Makes the singleton family Z coproduct X' -> X a ph covering.
\item Zariski sheaves send disjoint unions to products. Same chapter, the sheaf condition for the Zariski topology used at proof lines 10898-10902. Allows the expanded fibre product to be evaluated on F.
\item The blow up square and its isomorphism away from the exceptional divisor. Same chapter, equation labelled equation-blow-up-square, invoked in the statement. Supplies b : X' -> X and the isomorphism on the complements used to build the second covering.
\end{itemize}\medskip
\ifdefined\hypertarget\hypertarget{sourcepointer-1}{}\fi
\noindent\textit{Source reference.} Original definition 0DBL: sources/stacks-chow/tags.txt, source lines 16701--16701. Cross-chapter standard prerequisite located in the project tag list staged in the local snapshot; invoked at proof line None of the counted proof. It is retained only as a navigation pointer; every load-bearing step is inside the delivered span.\par\medskip
\ifdefined\hypertarget\hypertarget{sourcepointer-2}{}\fi
\noindent\textit{Source reference.} Original equation 0EV5: sources/stacks-flat/pinned-original-source.tex, source lines 10867--10867. Same-chapter pointer into the author environment equation at source line 10867; invoked at proof line None of the counted proof. It is retained only as a navigation pointer; every load-bearing step is inside the delivered span.\par\medskip
\ifdefined\hypertarget\hypertarget{sourcepointer-3}{}\fi
\noindent\textit{Source reference.} Original lemma 0DES: sources/stacks-chow/tags.txt, source lines 16812--16812. Cross-chapter standard prerequisite located in the project tag list staged in the local snapshot; invoked at proof line 10895 of the counted proof. It is retained only as a navigation pointer; every load-bearing step is inside the delivered span.\par\medskip
\begin{lemma}
\label{lemma-blow-up-square-ph}
Let $\mathcal{F}$ be a sheaf on a site $(\Sch/S)_{ph}$, see
Topologies, Definition \ref{topologies-definition-big-small-ph}.
Then for any blow up square (\ref{equation-blow-up-square})
in the category $(\Sch/S)_{ph}$ the diagram
$$
\xymatrix{
\mathcal{F}(E) & \mathcal{F}(X') \ar[l] \\
\mathcal{F}(Z) \ar[u] & \mathcal{F}(X) \ar[u] \ar[l]
}
$$
is cartesian in the category of sets.
\end{lemma}

\begin{proof}
Since $Z \amalg X' \to X$ is a surjective proper morphism
we see that $\{Z \amalg X' \to X\}$ is a ph covering
(Topologies, Lemma \ref{topologies-lemma-surjective-proper-ph}).
We have
$$
(Z \amalg X') \times_X (Z \amalg X') =
Z \amalg E \amalg E \amalg X' \times_X X'
$$
Since $\mathcal{F}$ is a Zariski sheaf we see that
$\mathcal{F}$ sends disjoint unions to products.
Thus the sheaf condition for the covering
$\{Z \amalg X' \to X\}$ says that
$\mathcal{F}(X) \to \mathcal{F}(Z) \times \mathcal{F}(X')$
is injective with image the set of pairs $(t, s')$ such that
(a) $t|_E = s'|_E$ and (b) $s'$ is in the equalizer of the two maps
$\mathcal{F}(X') \to \mathcal{F}(X' \times_X X')$.
Next, observe that the obvious morphism
$$
E \times_Z E \amalg X' \longrightarrow X' \times_X X'
$$
is a surjective proper morphism as $b$ induces
an isomorphism $X' \setminus E \to X \setminus Z$. We conclude that
$\mathcal{F}(X' \times_X X') \to
\mathcal{F}(E \times_Z E) \times \mathcal{F}(X')$ is injective.
It follows that (a) $\Rightarrow$ (b) which means that the lemma is true.
\end{proof}

\end{document}
